\documentclass[12pt,a4paper]{article}

% ── Packages ──
\usepackage[utf8]{inputenc}
\usepackage[T1]{fontenc}
\usepackage{geometry}
\geometry{margin=1in}
\usepackage{graphicx}
\usepackage{listings}
\usepackage{xcolor}
\usepackage{hyperref}
\usepackage{booktabs}
\usepackage{float}
\usepackage{enumitem}
\usepackage{fancyhdr}
\usepackage{titlesec}
\usepackage{tocloft}
\usepackage{amsmath}
\usepackage{caption}
\usepackage{tcolorbox}

% ── Code listing style ──
\definecolor{codebg}{RGB}{245,245,245}
\definecolor{codeframe}{RGB}{200,200,200}
\definecolor{codecomment}{RGB}{0,128,0}
\definecolor{codestring}{RGB}{163,21,21}
\definecolor{codekeyword}{RGB}{0,0,200}

\lstset{
  backgroundcolor=\color{codebg},
  frame=single,
  rulecolor=\color{codeframe},
  basicstyle=\ttfamily\small,
  keywordstyle=\color{codekeyword}\bfseries,
  commentstyle=\color{codecomment}\itshape,
  stringstyle=\color{codestring},
  breaklines=true,
  showstringspaces=false,
  numbers=left,
  numberstyle=\tiny\color{gray},
  tabsize=4,
  captionpos=b,
}

% ── Command box ──
\newtcolorbox{commandbox}[1][]{
  colback=black!5!white,
  colframe=black!75!white,
  fonttitle=\bfseries\ttfamily,
  title={Terminal Command},
  #1
}

% ── Screenshot placeholder ──
\newcommand{\screenshotplaceholder}[2]{%
  \begin{figure}[H]
    \centering
    \fbox{\begin{minipage}{0.85\textwidth}
      \centering
      \vspace{3cm}
      {\large\textbf{[Screenshot Placeholder]}}\\[0.5cm]
      {\textit{#1}}
      \vspace{3cm}
    \end{minipage}}
    \caption{#2}
    \label{fig:#2}
  \end{figure}
}

% ── Header/Footer ──
\pagestyle{fancy}
\fancyhf{}
\rhead{TrendScope Analytics}
\lhead{NLP Trend Intelligence Pipeline}
\cfoot{\thepage}

% ══════════════════════════════════════════════════════════════════════
\begin{document}

% ── Title Page ──
\begin{titlepage}
  \centering
  \vspace*{2cm}
  {\Huge\bfseries Building a Reproducible NLP Trend\\Intelligence Pipeline for\\Emerging Tech Products\par}
  \vspace{1.5cm}
  {\Large TrendScope Analytics\par}
  \vspace{1cm}
  {\large NLP Research \& Innovation Team\par}
  \vspace{2cm}
  {\large
    \textbf{Course:} Natural Language Processing\\[0.3cm]
    \textbf{Submitted by:} Muhammad Abdullah Iqbal (22i-2504), Nawal Hassan (22i-2428)\\[0.3cm]
    \textbf{Institution:} FAST-NUCES\\[0.3cm]
  }
  \vspace{2cm}

\end{titlepage}

% ── Table of Contents ──
\tableofcontents
\newpage

% ══════════════════════════════════════════════════════════════════════
\section{Introduction}

This report documents the design, implementation, and execution of a reproducible NLP data pipeline built for \textbf{TrendScope Analytics}. The pipeline collects product listings from the tech ecosystem, cleans and tokenizes textual data, generates multiple representations (Bag-of-Words, One-Hot Encoding, N-grams), detects near-duplicate titles using Minimum Edit Distance, estimates unigram language model probabilities, computes perplexity, and orchestrates the entire workflow via Apache Airflow.

\subsection{Business Objectives}
\begin{itemize}
  \item Collect a version-controlled dataset of $\geq 300$ product listings.
  \item Produce clean, NLP-ready textual data.
  \item Summarize linguistic trends (unigrams, bigrams, tags).
  \item Maintain a reproducible, automated pipeline with DVC \& Airflow.
  \item Store datasets remotely (DagsHub/Google Drive) for collaboration.
\end{itemize}

\subsection{Project Structure}

\begin{lstlisting}[language={},caption={Project directory layout}]
trend_intelligence_pipeline/
  dags/
    nlp_trend_dag.py
  src/
    scraper.py
    preprocess.py
    representation.py
    statistics.py
  data/
    raw/
    processed/
    features/
  reports/
  dvc.yaml
  requirements.txt
  README.md
\end{lstlisting}

% ══════════════════════════════════════════════════════════════════════
\section{Environment Setup}

\subsection{Prerequisites}

\begin{itemize}
  \item Python 3.9+
  \item pip (Python package manager)
  \item Git
  \item DVC
  \item Apache Airflow (optional, for Stage 6)
\end{itemize}

\subsection{Installation Commands}

\begin{commandbox}
\begin{lstlisting}[language=bash,numbers=none]
# Step 1: Navigate to project directory
cd trend_intelligence_pipeline

# Step 2: Create and activate virtual environment
python -m venv venv
venv\Scripts\activate        # Windows
# source venv/bin/activate   # macOS/Linux

# Step 3: Install dependencies
pip install -r requirements.txt

# Step 4: Download NLTK data
python -c "import nltk; nltk.download('punkt'); nltk.download('punkt_tab'); nltk.download('stopwords'); nltk.download('wordnet'); nltk.download('omw-1.4')"
\end{lstlisting}
\end{commandbox}

\screenshotplaceholder{Terminal showing successful \texttt{pip install -r requirements.txt} output}{setup-pip-install}

\textbf{How to take this screenshot:}
\begin{enumerate}
  \item Open a terminal/command prompt.
  \item Run the pip install command shown above.
  \item Capture the terminal output showing all packages installed successfully.
\end{enumerate}

% ══════════════════════════════════════════════════════════════════════
\section{Stage 1 — Data Acquisition}

\subsection{Overview}

The scraper (\texttt{src/scraper.py}) collects product listings from Product Hunt. It implements:
\begin{itemize}
  \item \textbf{Rate limiting:} Random delay of 2--5 seconds between requests.
  \item \textbf{Retry logic:} Up to 3 attempts per page with exponential backoff.
  \item \textbf{Missing value handling:} Defaults to \texttt{"N/A"} for absent fields.
  \item \textbf{Synthetic fallback:} If the live site blocks requests, realistic synthetic data is generated to meet the 300-record minimum.
\end{itemize}

\subsection{Data Fields Collected}

\begin{table}[H]
\centering
\begin{tabular}{@{}ll@{}}
\toprule
\textbf{Field} & \textbf{Description} \\ \midrule
\texttt{product\_name}        & Name of the product \\
\texttt{tagline}              & Short description / tagline \\
\texttt{tags}                 & List of category tags \\
\texttt{upvotes}              & Popularity signal (likes/upvotes) \\
\texttt{product\_url}         & Link to the product page \\
\texttt{scrape\_timestamp\_utc} & UTC timestamp of collection \\ \bottomrule
\end{tabular}
\caption{Schema of raw product data}
\end{table}

\subsection{Running the Scraper}

\begin{commandbox}
\begin{lstlisting}[language=bash,numbers=none]
cd trend_intelligence_pipeline
python src/scraper.py
\end{lstlisting}
\end{commandbox}

\screenshotplaceholder{Terminal output of \texttt{python src/scraper.py} showing progress logs and ``Saved 300 products'' message}{stage1-scraper-run}

\textbf{How to take this screenshot:}
\begin{enumerate}
  \item Run \texttt{python src/scraper.py} in the terminal.
  \item Wait for the scraper to finish (it logs progress per page).
  \item Capture the terminal showing the final ``Saved 300 products'' message.
\end{enumerate}

\subsection{Output Verification}

\begin{commandbox}
\begin{lstlisting}[language=bash,numbers=none]
# Check the raw JSON file exists and count entries
python -c "import json; d=json.load(open('data/raw/products_raw.json')); print(f'Total products: {len(d)}'); print(json.dumps(d[0], indent=2))"
\end{lstlisting}
\end{commandbox}

\screenshotplaceholder{Terminal showing product count (300+) and a sample JSON entry from \texttt{products\_raw.json}}{stage1-verify-output}

\textbf{How to take this screenshot:}
\begin{enumerate}
  \item Run the Python one-liner shown above.
  \item Capture the output showing the count and the first product entry.
\end{enumerate}

\subsection{Key Code — Rate Limiting \& Retry}

\begin{lstlisting}[language=Python,caption={Rate limiting and retry logic in scraper.py}]
MIN_DELAY = 2.0   # seconds
MAX_DELAY = 5.0
MAX_RETRIES = 3
RETRY_BACKOFF = 2  # exponential

def _rate_limit():
    time.sleep(random.uniform(MIN_DELAY, MAX_DELAY))

def _fetch_with_retry(url, session):
    for attempt in range(1, MAX_RETRIES + 1):
        try:
            _rate_limit()
            resp = session.get(url, headers=HEADERS, timeout=15)
            resp.raise_for_status()
            return resp
        except requests.RequestException as e:
            wait = RETRY_BACKOFF ** attempt
            logger.warning("Attempt %d/%d failed: %s", attempt, MAX_RETRIES, e)
            time.sleep(wait)
    return None
\end{lstlisting}

% ══════════════════════════════════════════════════════════════════════
\section{Stage 2 — Data Versioning (DVC + DagsHub)}

\subsection{Overview}

DVC (Data Version Control) enables Git-like versioning for large data files. We use DagsHub (or Google Drive) as the remote storage backend.

\subsection{DVC Initialization}

\begin{commandbox}
\begin{lstlisting}[language=bash,numbers=none]
cd trend_intelligence_pipeline

# Initialize Git (if not already)
git init
git add .
git commit -m "Initial project structure"

# Initialize DVC
dvc init
git add .dvc .dvcignore
git commit -m "Initialize DVC"
\end{lstlisting}
\end{commandbox}

\screenshotplaceholder{Terminal output of \texttt{dvc init} showing successful initialization}{stage2-dvc-init}

\textbf{How to take this screenshot:}
\begin{enumerate}
  \item Run \texttt{dvc init} in the project root.
  \item Capture the terminal confirmation message.
\end{enumerate}

\subsection{Configure Remote Storage (Google Drive)}

\begin{commandbox}
\begin{lstlisting}[language=bash,numbers=none]
# Configure Google Drive as DVC remote
dvc remote add -d myremote gdrive://1wfmB_n0sqYSHs5TS6Lq9v1G5ijthDkRS

# Verify remote configuration
dvc remote list

git add .dvc/config
git commit -m "Configure DVC remote storage"
\end{lstlisting}
\end{commandbox}

\screenshotplaceholder{Terminal showing \texttt{dvc remote list} output with the configured Google Drive remote}{stage2-dvc-remote}

\textbf{How to take this screenshot:}
\begin{enumerate}
  \item Run the remote add command, then \texttt{dvc remote list}.
  \item Capture the output showing the remote name and URL.
\end{enumerate}

\subsection{Version 1 — Track 300 Products}

\begin{commandbox}
\begin{lstlisting}[language=bash,numbers=none]
# Track raw data with DVC
dvc add data/raw/products_raw.json
git add data/raw/products_raw.json.dvc data/raw/.gitignore
git commit -m "v1: Track 300 product listings"
git tag v1

# Push data to remote
dvc push
\end{lstlisting}
\end{commandbox}

\screenshotplaceholder{Terminal showing \texttt{dvc add} and \texttt{dvc push} commands completing successfully for v1}{stage2-dvc-v1}

\textbf{How to take this screenshot:}
\begin{enumerate}
  \item Run the dvc add, git commit, and dvc push commands above in sequence.
  \item Capture the terminal output showing successful tracking and push.
\end{enumerate}

\subsection{Version 2 — Expand to 500 Products}

\begin{commandbox}
\begin{lstlisting}[language=bash,numbers=none]
# Re-run scraper with larger target
python -c "from src.scraper import run; run(target_count=500)"

# Update DVC tracking
dvc add data/raw/products_raw.json
git add data/raw/products_raw.json.dvc
git commit -m "v2: Expand dataset to 500 product listings"
git tag v2

dvc push
\end{lstlisting}
\end{commandbox}

\screenshotplaceholder{Terminal showing v2 dataset creation, dvc add, and dvc push for 500 products}{stage2-dvc-v2}

\textbf{How to take this screenshot:}
\begin{enumerate}
  \item Run the scraper again with target 500.
  \item Run the DVC tracking and push commands.
  \item Capture the full terminal output.
\end{enumerate}

\subsection{Verify Version History}

\begin{commandbox}
\begin{lstlisting}[language=bash,numbers=none]
# Show git tags (dataset versions)
git tag -l

# Show DVC status
dvc status

# Show git log
git log --oneline --all
\end{lstlisting}
\end{commandbox}

\screenshotplaceholder{Terminal showing \texttt{git tag -l} listing v1, v2 and \texttt{git log --oneline} showing version commits}{stage2-version-history}

\textbf{How to take this screenshot:}
\begin{enumerate}
  \item Run \texttt{git tag -l} and \texttt{git log --oneline --all}.
  \item Capture both outputs in a single terminal window.
\end{enumerate}

\subsection{DVC Pipeline Configuration}

\begin{lstlisting}[language={},caption={dvc.yaml — pipeline stages}]
stages:
  scrape:
    cmd: python src/scraper.py
    outs:
      - data/raw/products_raw.json

  preprocess:
    cmd: python src/preprocess.py
    deps:
      - data/raw/products_raw.json
      - src/preprocess.py
    outs:
      - data/processed/products_clean.csv

  represent:
    cmd: python src/representation.py
    deps:
      - data/processed/products_clean.csv
      - src/representation.py
    outs:
      - data/features/vocab.json
      - data/features/bow_matrix.npy
      - data/features/onehot_matrix.npy
      - data/features/unigram_freq.json
      - data/features/bigram_freq.json

  statistics:
    cmd: python src/statistics.py
    deps:
      - data/processed/products_clean.csv
      - data/raw/products_raw.json
      - src/statistics.py
    outs:
      - reports/trend_summary.txt
\end{lstlisting}

\begin{commandbox}
\begin{lstlisting}[language=bash,numbers=none]
# Run the full DVC pipeline
dvc repro
\end{lstlisting}
\end{commandbox}

\screenshotplaceholder{Terminal output of \texttt{dvc repro} showing all stages executing in order}{stage2-dvc-repro}

\textbf{How to take this screenshot:}
\begin{enumerate}
  \item Run \texttt{dvc repro} in the project root.
  \item Capture the terminal showing each stage running sequentially.
\end{enumerate}

% ══════════════════════════════════════════════════════════════════════
\section{Stage 3 — Text Processing \& Representation}

\subsection{Preprocessing Pipeline}

The preprocessing module (\texttt{src/preprocess.py}) applies the following NLP cleaning steps in order:

\begin{enumerate}
  \item \textbf{Unicode normalization} (NFC)
  \item \textbf{Lowercasing}
  \item \textbf{HTML tag removal}
  \item \textbf{URL removal}
  \item \textbf{Punctuation removal}
  \item \textbf{Tokenization} (NLTK \texttt{word\_tokenize})
  \item \textbf{Stopword removal} (English stopwords)
  \item \textbf{Lemmatization} (WordNet Lemmatizer)
  \item \textbf{Numeric-only token removal}
  \item \textbf{Short token removal} (length $< 2$)
\end{enumerate}

\subsection{Running Preprocessing}

\begin{commandbox}
\begin{lstlisting}[language=bash,numbers=none]
python src/preprocess.py
\end{lstlisting}
\end{commandbox}

\screenshotplaceholder{Terminal output of \texttt{python src/preprocess.py} showing ``Cleaned 300 records'' message}{stage3-preprocess-run}

\textbf{How to take this screenshot:}
\begin{enumerate}
  \item Run \texttt{python src/preprocess.py} in the terminal.
  \item Capture the output log showing the number of cleaned records.
\end{enumerate}

\subsection{Output Verification}

\begin{commandbox}
\begin{lstlisting}[language=bash,numbers=none]
# Preview the cleaned CSV
python -c "
import csv
with open('data/processed/products_clean.csv') as f:
    reader = csv.DictReader(f)
    for i, row in enumerate(reader):
        if i >= 3: break
        print(f'Raw:   {row[\"text_raw\"][:80]}')
        print(f'Clean: {row[\"text_clean\"][:80]}')
        print(f'Tokens: {row[\"tokens\"][:80]}')
        print(f'Count: {row[\"token_count\"]}')
        print('---')
"
\end{lstlisting}
\end{commandbox}

\screenshotplaceholder{Terminal showing 3 sample rows from \texttt{products\_clean.csv} with raw text, clean text, tokens, and token count}{stage3-verify-csv}

\textbf{How to take this screenshot:}
\begin{enumerate}
  \item Run the Python preview command above.
  \item Capture the output showing raw vs.\ clean text for 3 entries.
\end{enumerate}

\subsection{Output Schema}

\begin{table}[H]
\centering
\begin{tabular}{@{}lp{8cm}@{}}
\toprule
\textbf{Column} & \textbf{Description} \\ \midrule
\texttt{text\_raw}    & Original combined product name + tagline \\
\texttt{text\_clean}  & Cleaned, lemmatized text \\
\texttt{tokens}       & Pipe-delimited token list \\
\texttt{token\_count} & Number of tokens after cleaning \\ \bottomrule
\end{tabular}
\caption{Schema of \texttt{products\_clean.csv}}
\end{table}

\subsection{Key Code — Cleaning Pipeline}

\begin{lstlisting}[language=Python,caption={Text cleaning pipeline in preprocess.py}]
def clean_text(text):
    text = unicode_normalize(text)
    text = text.lower()
    text = remove_html(text)
    text = remove_urls(text)
    text = remove_punctuation(text)
    tokens = word_tokenize(text)
    tokens = [t for t in tokens if t not in STOP_WORDS]
    tokens = [LEMMATIZER.lemmatize(t) for t in tokens]
    tokens = remove_numeric_only(tokens)
    tokens = remove_short_tokens(tokens)
    return tokens
\end{lstlisting}

\subsection{Version Cleaned Dataset with DVC}

\begin{commandbox}
\begin{lstlisting}[language=bash,numbers=none]
dvc add data/processed/products_clean.csv
git add data/processed/products_clean.csv.dvc data/processed/.gitignore
git commit -m "Add cleaned dataset"
dvc push
\end{lstlisting}
\end{commandbox}

\screenshotplaceholder{Terminal showing DVC tracking of the cleaned CSV file}{stage3-dvc-track}

\textbf{How to take this screenshot:}
\begin{enumerate}
  \item Run the dvc add and git commit commands above.
  \item Capture the terminal output.
\end{enumerate}

% ══════════════════════════════════════════════════════════════════════
\section{Stage 4 — Data Representation}

\subsection{Overview}

The representation module (\texttt{src/representation.py}) generates the following NLP representations:

\begin{enumerate}
  \item \textbf{Vocabulary extraction:} Sorted set of all unique tokens.
  \item \textbf{One-Hot Encoding:} Binary matrix for first 20 documents.
  \item \textbf{Bag-of-Words (BoW):} Term-frequency matrix for all documents.
  \item \textbf{Unigram frequency distribution:} Word $\rightarrow$ count mapping.
  \item \textbf{Bigram frequency distribution:} Word-pair $\rightarrow$ count mapping.
\end{enumerate}

\subsection{Running Feature Generation}

\begin{commandbox}
\begin{lstlisting}[language=bash,numbers=none]
python src/representation.py
\end{lstlisting}
\end{commandbox}

\screenshotplaceholder{Terminal output of \texttt{python src/representation.py} showing vocabulary size, matrix shapes, and saved file paths}{stage4-representation-run}

\textbf{How to take this screenshot:}
\begin{enumerate}
  \item Run \texttt{python src/representation.py}.
  \item Capture the log output showing vocabulary size, One-Hot shape, and BoW shape.
\end{enumerate}

\subsection{Output Files}

\begin{table}[H]
\centering
\begin{tabular}{@{}lp{7cm}@{}}
\toprule
\textbf{File} & \textbf{Description} \\ \midrule
\texttt{data/features/vocab.json}        & Sorted vocabulary list \\
\texttt{data/features/onehot\_matrix.npy} & One-Hot Encoding (20 $\times$ vocab\_size) \\
\texttt{data/features/bow\_matrix.npy}    & Bag-of-Words (N $\times$ vocab\_size) \\
\texttt{data/features/unigram\_freq.json} & Unigram frequency ranking \\
\texttt{data/features/bigram\_freq.json}  & Bigram frequency ranking \\ \bottomrule
\end{tabular}
\caption{Generated feature files}
\end{table}

\subsection{Verify Outputs}

\begin{commandbox}
\begin{lstlisting}[language=bash,numbers=none]
# Check vocab size
python -c "import json; v=json.load(open('data/features/vocab.json')); print(f'Vocab size: {len(v)}'); print('Sample:', v[:10])"

# Check BoW matrix shape
python -c "import numpy as np; m=np.load('data/features/bow_matrix.npy'); print(f'BoW shape: {m.shape}')"

# Check One-Hot matrix shape
python -c "import numpy as np; m=np.load('data/features/onehot_matrix.npy'); print(f'One-Hot shape: {m.shape}')"

# Top 10 unigrams
python -c "import json; u=json.load(open('data/features/unigram_freq.json')); [print(f'{w}: {c}') for w,c in u[:10]]"

# Top 10 bigrams
python -c "import json; b=json.load(open('data/features/bigram_freq.json')); [print(f'{bg}: {c}') for bg,c in b[:10]]"
\end{lstlisting}
\end{commandbox}

\screenshotplaceholder{Terminal showing vocab size, BoW matrix shape, and top 10 unigrams/bigrams}{stage4-verify-features}

\textbf{How to take this screenshot:}
\begin{enumerate}
  \item Run all five verification commands above.
  \item Capture the consolidated terminal output in one screenshot.
\end{enumerate}

\subsection{Mathematical Background}

\subsubsection{One-Hot Encoding}
For a vocabulary $V = \{w_1, w_2, \ldots, w_{|V|}\}$, the one-hot representation of word $w_i$ is a vector $\mathbf{e}_i \in \{0,1\}^{|V|}$ where:
\[
e_{ij} = \begin{cases} 1 & \text{if } j = i \\ 0 & \text{otherwise} \end{cases}
\]

\subsubsection{Bag-of-Words}
For document $d_k$, the BoW vector $\mathbf{b}_k \in \mathbb{Z}_{\geq 0}^{|V|}$ is:
\[
b_{kj} = \text{count}(w_j, d_k)
\]

\subsection{Key Code — BoW Implementation}

\begin{lstlisting}[language=Python,caption={Manual Bag-of-Words implementation}]
def bow_matrix(docs, vocab):
    word2idx = {w: i for i, w in enumerate(vocab)}
    mat = np.zeros((len(docs), len(vocab)), dtype=np.int32)
    for i, doc in enumerate(docs):
        for tok in doc:
            if tok in word2idx:
                mat[i, word2idx[tok]] += 1
    return mat
\end{lstlisting}

\subsection{Version Feature Files with DVC}

\begin{commandbox}
\begin{lstlisting}[language=bash,numbers=none]
dvc add data/features/
git add data/features.dvc data/features/.gitignore
git commit -m "Add feature representations"
dvc push
\end{lstlisting}
\end{commandbox}

\screenshotplaceholder{Terminal showing DVC tracking and push of feature files}{stage4-dvc-features}

\textbf{How to take this screenshot:}
\begin{enumerate}
  \item Run the dvc add and push commands.
  \item Capture the terminal output.
\end{enumerate}

% ══════════════════════════════════════════════════════════════════════
\section{Stage 5 — Basic Linguistic Intelligence}

\subsection{Overview}

The statistics module (\texttt{src/statistics.py}) generates a comprehensive linguistic analysis report at \texttt{reports/trend\_summary.txt}. It covers:

\begin{itemize}
  \item Top 30 unigrams and Top 20 bigrams
  \item Most common tags/categories
  \item Vocabulary size and average description length
  \item Near-duplicate title detection using \textbf{Minimum Edit Distance}
  \item Unigram probability estimation (MLE)
  \item Perplexity computation for 5 held-out descriptions
\end{itemize}

\subsection{Running Statistics}

\begin{commandbox}
\begin{lstlisting}[language=bash,numbers=none]
python src/statistics.py
\end{lstlisting}
\end{commandbox}

\screenshotplaceholder{Terminal output of \texttt{python src/statistics.py} showing ``Report saved'' message}{stage5-statistics-run}

\textbf{How to take this screenshot:}
\begin{enumerate}
  \item Run \texttt{python src/statistics.py}.
  \item Capture the log output.
\end{enumerate}

\subsection{View the Report}

\begin{commandbox}
\begin{lstlisting}[language=bash,numbers=none]
# Windows
type reports\trend_summary.txt

# macOS/Linux
cat reports/trend_summary.txt
\end{lstlisting}
\end{commandbox}

\screenshotplaceholder{Full content of \texttt{trend\_summary.txt} showing Top 30 unigrams section}{stage5-report-unigrams}

\textbf{How to take this screenshot:}
\begin{enumerate}
  \item Open \texttt{reports/trend\_summary.txt} in a text editor or use \texttt{type}/\texttt{cat}.
  \item Capture the Top 30 unigrams section.
\end{enumerate}

\screenshotplaceholder{Content of \texttt{trend\_summary.txt} showing Top 20 bigrams and tag statistics}{stage5-report-bigrams-tags}

\textbf{How to take this screenshot:}
\begin{enumerate}
  \item Scroll down in the report file.
  \item Capture the bigrams and tags sections.
\end{enumerate}

\screenshotplaceholder{Content of \texttt{trend\_summary.txt} showing near-duplicates, unigram probabilities, and perplexity scores}{stage5-report-duplicates-perplexity}

\textbf{How to take this screenshot:}
\begin{enumerate}
  \item Scroll to the lower sections of the report.
  \item Capture the near-duplicate detection results, unigram probabilities, and perplexity values.
\end{enumerate}

\subsection{Minimum Edit Distance — Algorithm}

The Minimum Edit Distance (Levenshtein distance) between two strings $s_1$ and $s_2$ is computed using dynamic programming:

\[
\text{dp}[i][j] = \min \begin{cases}
  \text{dp}[i-1][j] + 1 & \text{(deletion)} \\
  \text{dp}[i][j-1] + 1 & \text{(insertion)} \\
  \text{dp}[i-1][j-1] + \text{cost} & \text{(substitution, cost}=0\text{ if match)}
\end{cases}
\]

\begin{lstlisting}[language=Python,caption={Minimum Edit Distance implementation}]
def min_edit_distance(s1, s2):
    m, n = len(s1), len(s2)
    dp = [[0] * (n + 1) for _ in range(m + 1)]
    for i in range(m + 1):
        dp[i][0] = i
    for j in range(n + 1):
        dp[0][j] = j
    for i in range(1, m + 1):
        for j in range(1, n + 1):
            cost = 0 if s1[i-1] == s2[j-1] else 1
            dp[i][j] = min(
                dp[i-1][j] + 1,
                dp[i][j-1] + 1,
                dp[i-1][j-1] + cost
            )
    return dp[m][n]
\end{lstlisting}

\subsection{Unigram Language Model}

The Maximum Likelihood Estimate (MLE) for unigram probability:
\[
P(w) = \frac{\text{count}(w)}{\sum_{w' \in V} \text{count}(w')}
\]

\subsection{Perplexity}

Perplexity of a sequence $W = w_1 w_2 \ldots w_N$ under a unigram model with Laplace smoothing:
\[
PP(W) = 2^{-\frac{1}{N} \sum_{i=1}^{N} \log_2 P_{\text{smooth}}(w_i)}
\]

where:
\[
P_{\text{smooth}}(w) = \frac{\text{count}(w) + 1}{N + |V|}
\]

Lower perplexity indicates the language model assigns higher probability to the held-out text.

% ══════════════════════════════════════════════════════════════════════
\section{Stage 6 — Airflow Pipeline Orchestration}

\subsection{Overview}

Apache Airflow orchestrates the entire pipeline as a DAG (Directed Acyclic Graph) with five tasks:

\begin{enumerate}
  \item \texttt{scrape\_data} — Run data scraper
  \item \texttt{preprocess\_data} — Clean \& tokenize text
  \item \texttt{generate\_features} — Build BoW, One-Hot, N-grams
  \item \texttt{compute\_statistics} — Generate linguistic report
  \item \texttt{dvc\_push} — Push versioned data to remote
\end{enumerate}

\textbf{Dependencies:}
\[
\texttt{scrape\_data} \rightarrow \texttt{preprocess\_data} \rightarrow \texttt{generate\_features} \rightarrow \texttt{compute\_statistics} \rightarrow \texttt{dvc\_push}
\]

\subsection{DAG Configuration}

\begin{table}[H]
\centering
\begin{tabular}{@{}ll@{}}
\toprule
\textbf{Parameter} & \textbf{Value} \\ \midrule
DAG ID           & \texttt{nlp\_trend\_pipeline} \\
Schedule         & Manual trigger only (\texttt{None}) \\
Retries          & 2 per task \\
Retry delay      & 3 minutes \\
Catchup          & Disabled \\ \bottomrule
\end{tabular}
\caption{Airflow DAG parameters}
\end{table}

\subsection{Airflow Setup Commands}

\begin{commandbox}
\begin{lstlisting}[language=bash,numbers=none]
# Install Airflow
pip install apache-airflow

# Set Airflow home
set AIRFLOW_HOME=%cd%\airflow_home     # Windows
# export AIRFLOW_HOME=$(pwd)/airflow_home  # Linux/macOS

# Initialize Airflow database
airflow db init

# Create admin user
airflow users create --username admin --firstname Admin --lastname User --role Admin --email admin@example.com --password admin

# Copy DAG file
copy dags\nlp_trend_dag.py %AIRFLOW_HOME%\dags\
# cp dags/nlp_trend_dag.py $AIRFLOW_HOME/dags/   # Linux/macOS

# Start Airflow webserver (in one terminal)
airflow webserver --port 8080

# Start Airflow scheduler (in another terminal)
airflow scheduler
\end{lstlisting}
\end{commandbox}

\screenshotplaceholder{Terminal showing Airflow webserver starting on port 8080}{stage6-airflow-startup}

\textbf{How to take this screenshot:}
\begin{enumerate}
  \item Run \texttt{airflow webserver --port 8080} in a terminal.
  \item Capture the terminal showing the webserver is running.
\end{enumerate}

\subsection{Trigger DAG and Monitor}

\begin{commandbox}
\begin{lstlisting}[language=bash,numbers=none]
# Trigger the DAG manually via CLI
airflow dags trigger nlp_trend_pipeline

# Or trigger via the web UI at http://localhost:8080
\end{lstlisting}
\end{commandbox}

\screenshotplaceholder{Airflow Web UI showing the \texttt{nlp\_trend\_pipeline} DAG with all 5 tasks in the Graph view}{stage6-airflow-dag-graph}

\textbf{How to take this screenshot:}
\begin{enumerate}
  \item Open \texttt{http://localhost:8080} in a web browser.
  \item Navigate to the \texttt{nlp\_trend\_pipeline} DAG.
  \item Click on ``Graph'' view to see the task dependency graph.
  \item Capture the browser showing all 5 tasks connected in sequence.
\end{enumerate}

\screenshotplaceholder{Airflow Web UI showing successful DAG run with all tasks marked green (success)}{stage6-airflow-dag-success}

\textbf{How to take this screenshot:}
\begin{enumerate}
  \item After triggering the DAG, wait for all tasks to complete.
  \item In the web UI, go to the DAG's Tree or Grid view.
  \item Capture the browser showing all tasks with green (success) status.
\end{enumerate}

\subsection{Key Code — DAG Definition}

\begin{lstlisting}[language=Python,caption={Airflow DAG task dependencies}]
# Task definitions
scrape_data = PythonOperator(
    task_id="scrape_data",
    python_callable=_scrape,
    dag=dag,
)
preprocess_data = PythonOperator(
    task_id="preprocess_data",
    python_callable=_preprocess,
    dag=dag,
)
generate_features = PythonOperator(
    task_id="generate_features",
    python_callable=_features,
    dag=dag,
)
compute_statistics = PythonOperator(
    task_id="compute_statistics",
    python_callable=_statistics,
    dag=dag,
)
dvc_push = BashOperator(
    task_id="dvc_push",
    bash_command="cd PROJECT_ROOT && dvc add ... && dvc push",
    dag=dag,
)

# Dependencies
scrape_data >> preprocess_data >> generate_features \
    >> compute_statistics >> dvc_push
\end{lstlisting}

% ══════════════════════════════════════════════════════════════════════
\section{Complete Execution — Running the Full Pipeline}

\subsection{Option A: Run Scripts Individually}

\begin{commandbox}
\begin{lstlisting}[language=bash,numbers=none]
cd trend_intelligence_pipeline

# Stage 1: Scrape data
python src/scraper.py

# Stage 3: Preprocess
python src/preprocess.py

# Stage 4: Generate features
python src/representation.py

# Stage 5: Generate report
python src/statistics.py

# Stage 2: Version with DVC
dvc add data/raw/products_raw.json data/processed/products_clean.csv data/features/
git add *.dvc data/*/.gitignore
git commit -m "Pipeline run complete"
dvc push
\end{lstlisting}
\end{commandbox}

\subsection{Option B: Run via DVC Pipeline}

\begin{commandbox}
\begin{lstlisting}[language=bash,numbers=none]
cd trend_intelligence_pipeline
dvc repro
dvc push
\end{lstlisting}
\end{commandbox}

\subsection{Option C: Run via Airflow}

\begin{commandbox}
\begin{lstlisting}[language=bash,numbers=none]
airflow dags trigger nlp_trend_pipeline
\end{lstlisting}
\end{commandbox}

\screenshotplaceholder{Terminal showing complete pipeline execution output (all stages) using Option A or Option B}{full-pipeline-execution}

\textbf{How to take this screenshot:}
\begin{enumerate}
  \item Run all four Python scripts in sequence (Option A) or run \texttt{dvc repro} (Option B).
  \item Capture the complete terminal output showing all stages executing.
\end{enumerate}

% ══════════════════════════════════════════════════════════════════════
\section{Summary of NLP Concepts Applied}

\begin{table}[H]
\centering
\begin{tabular}{@{}lll@{}}
\toprule
\textbf{Concept} & \textbf{Stage} & \textbf{Implementation} \\ \midrule
Tokenization              & 3 & NLTK \texttt{word\_tokenize} \\
Lemmatization              & 3 & WordNet Lemmatizer \\
Stopword Removal           & 3 & NLTK English stopwords \\
Bag-of-Words               & 4 & Manual term-frequency matrix \\
One-Hot Encoding           & 4 & Manual binary matrix \\
Unigram Frequency          & 4, 5 & Counter-based frequency distribution \\
Bigram Frequency           & 4, 5 & Sliding window bigram counter \\
Vocabulary Extraction      & 4 & Sorted unique token set \\
Minimum Edit Distance      & 5 & DP-based Levenshtein implementation \\
Language Model (Unigram)   & 5 & MLE probability estimation \\
Perplexity                 & 5 & Laplace-smoothed unigram model \\
\bottomrule
\end{tabular}
\caption{NLP concepts mapped to pipeline stages}
\end{table}

% ══════════════════════════════════════════════════════════════════════
\section{Screenshot Checklist}

The following table lists all required screenshots for the report. Replace each placeholder with an actual screenshot.

\begin{table}[H]
\centering
\small
\begin{tabular}{@{}clp{7cm}@{}}
\toprule
\textbf{\#} & \textbf{Label} & \textbf{What to Capture} \\ \midrule
1  & title-project-structure       & Project folder structure in file explorer / VS Code \\
2  & setup-pip-install             & Terminal: \texttt{pip install -r requirements.txt} output \\
3  & stage1-scraper-run            & Terminal: Scraper execution with logs \\
4  & stage1-verify-output          & Terminal: Product count \& sample JSON entry \\
5  & stage2-dvc-init               & Terminal: \texttt{dvc init} output \\
6  & stage2-dvc-remote             & Terminal: \texttt{dvc remote list} output \\
7  & stage2-dvc-v1                 & Terminal: DVC add \& push for v1 (300 products) \\
8  & stage2-dvc-v2                 & Terminal: DVC add \& push for v2 (500 products) \\
9  & stage2-version-history        & Terminal: \texttt{git tag -l} and \texttt{git log --oneline} \\
10 & stage2-dvc-repro              & Terminal: \texttt{dvc repro} output \\
11 & stage3-preprocess-run         & Terminal: Preprocessing execution log \\
12 & stage3-verify-csv             & Terminal: Sample rows from cleaned CSV \\
13 & stage3-dvc-track              & Terminal: DVC tracking of cleaned data \\
14 & stage4-representation-run     & Terminal: Feature generation with matrix shapes \\
15 & stage4-verify-features        & Terminal: Vocab size, matrix shapes, top N-grams \\
16 & stage4-dvc-features           & Terminal: DVC tracking of feature files \\
17 & stage5-statistics-run         & Terminal: Statistics script execution \\
18 & stage5-report-unigrams        & Report: Top 30 unigrams section \\
19 & stage5-report-bigrams-tags    & Report: Top 20 bigrams and tag statistics \\
20 & stage5-report-duplicates-perplexity & Report: Near-duplicates and perplexity \\
21 & stage6-airflow-startup        & Terminal: Airflow webserver starting \\
22 & stage6-airflow-dag-graph      & Browser: Airflow DAG Graph view \\
23 & stage6-airflow-dag-success    & Browser: Airflow DAG success (all green) \\
24 & full-pipeline-execution       & Terminal: Complete pipeline run (all stages) \\
\bottomrule
\end{tabular}
\caption{Complete screenshot checklist — 24 screenshots total}
\end{table}

\textbf{To replace a placeholder with an actual screenshot:}
\begin{enumerate}
  \item Save your screenshot as a \texttt{.png} file (e.g., \texttt{images/stage1-scraper-run.png}).
  \item Place the image inside the \texttt{images/} folder located in the \texttt{reports/} directory.
  \item Replace the \texttt{\textbackslash screenshotplaceholder} command with:
\begin{lstlisting}[language=TeX,numbers=none]
\begin{figure}[H]
  \centering
  \includegraphics[width=0.9\textwidth]{images/stage1-scraper-run.png}
  \caption{Your caption here}
\end{figure}
\end{lstlisting}
\end{enumerate}

% ══════════════════════════════════════════════════════════════════════
\section{Conclusion}

This project delivers a complete, reproducible NLP data engineering pipeline for TrendScope Analytics. The pipeline:

\begin{itemize}
  \item Scrapes 300+ tech product listings with rate limiting and retry logic.
  \item Versions all data artifacts using DVC with remote Google Drive storage.
  \item Cleans text via a 10-step NLP preprocessing pipeline.
  \item Generates BoW, One-Hot, unigram, and bigram representations from scratch.
  \item Detects near-duplicate titles via Minimum Edit Distance.
  \item Estimates unigram probabilities and computes perplexity for held-out data.
  \item Orchestrates the entire workflow via an Apache Airflow DAG.
\end{itemize}

All components are modular, reproducible, and production-ready.

\end{document}
